\section{Introduction}
\label{sec:introduction}

The Arnold fixed-point problem compares periodic orbits of Hamiltonian systems with critical points of smooth functions. We consider the proposed lower bound given by the stable Morse number, with stabilization understood geometrically.

\begin{definition}
Let $W$ be a connected closed smooth manifold. Its \emph{stable Morse number} $\SM(W)$ is the minimum of $\#\Crit(F)$ over all integers $k\geq0$ and smooth Morse functions
\[
 F:W\times\R^k\longrightarrow\R
\]
such that, when $k>0$, there are a nondegenerate quadratic form $Q$ on $\R^k$ and $R>0$ with $F(x,z)=Q(z)$ for every $x\in W$ and $|z|\geq R$. The signature of $Q$ is unrestricted. For $k=0$, there is no condition at infinity. The ordinary Morse number $\Morse(W)$ is the same minimum restricted to $k=0$.
\end{definition}

These functions have finitely many critical points, and $\SM(W)\leq\Morse(W)$. The definition is a minimum over actual smooth functions; an algebraic complex of small rank alone does not provide an upper bound.

Let $(M,\omega)$ be a connected closed symplectic manifold, and let $H\in C^\infty((\R/\Z)\times M,\R)$. We use the convention $\omega(X_{H_t},\cdot)=dH_t$ and denote its flow by $\phi_H^t$. A fixed point $x$ of $\phi_H^1$ is \emph{nondegenerate} if $1$ is not an eigenvalue of $d_x\phi_H^1$. Write
\[
 \Fix_0(\phi_H^1;H)=\{x\in\Fix(\phi_H^1):t\mapsto\phi_H^t(x)\text{ is contractible in }M\}.
\]
This counts distinct initial points, with no multiplicity from choices of capping disks. The stable Morse version of the fixed-point bound asks whether
\begin{equation}\label{eq:proposed-bound}
 \#\Fix_0(\phi_H^1;H)\geq\SM(M)
\end{equation}
holds whenever every fixed point is nondegenerate. The following theorem disproves this unrestricted assertion.

\begin{theorem}\label{thm:main}
There exist a simply connected closed Kähler manifold $(M,\omega)$ and a smooth one-periodic Hamiltonian $H$ such that every fixed point of $\phi_H^1$ is nondegenerate and
\[
 \#\Fix_0(\phi_H^1;H)=\#\Fix(\phi_H^1)=\SM(M)-16.
\]
One may take $\dim_{\R}M=1412$, with $\SM(M)=318968$ and $\#\Fix(\phi_H^1)=318952$.
\end{theorem}

\subsection*{Context and prior work}
Arnold's fixed-point questions ask how much of the critical-point theory of a small autonomous Hamiltonian persists for a general Hamiltonian path \cite{Arnold1986}. Conley and Zehnder established the torus case by variational methods \cite{ConleyZehnder1983}. Floer's use of pseudoholomorphic curves turned the problem into a comparison between a Hamiltonian chain complex and the topology of the ambient manifold \cite{Floer1989}. Hofer--Salamon and Ono extended the theory in the weakly monotone setting \cite{HoferSalamon1995,Ono1995}; virtual-moduli methods subsequently yielded the general rational homological bound through work of Fukaya--Ono, Liu--Tian, and Ruan \cite{FukayaOno1999,LiuTian1998,Ruan1999}. The general rational bound counts Betti numbers. The Morse and stable Morse formulations ask for a geometric critical-point minimum, which can be larger.

The stable formulation is recorded explicitly in \cite[Conjecture~1.2]{Golovko2020}. Its distinction from the ordinary Morse formulation is essential: Damian exhibited closed manifolds with different ordinary and stable Morse numbers \cite{Damian2002}. In the positive direction, Dimitroglou Rizell and Golovko proved the stable Morse lower bound for generic Hamiltonians when both the symplectic area class and the first Chern class vanish on $\pi_2(M)$ \cite[Corollary~1.4]{DRG2017}. Their argument controls the simple homotopy type of the Floer complex. An independent proof was obtained by P\"oder Balkest\aa hl \cite[Corollary~1.2]{PoderBalkestahl2021}. The almost-quadratic convention for stable Morse functions used by Dimitroglou Rizell and Golovko agrees with ours by Lemma~\ref{lem:infinity-conventions}. Theorem~\ref{thm:main} disproves the unrestricted stable formulation, including its version that counts all fixed points.

Recent integral refinements explain why torsion alone does not force the bound in \eqref{eq:proposed-bound}. Abouzaid and Blumberg proved homological lower bounds over fields of arbitrary positive characteristic \cite{AbouzaidBlumberg2021}. Bai and Xu obtained an integral bound incorporating torsion at all primes \cite[Theorem~A]{BaiXu2022}, and their subsequent foundations give the same bound in \cite[Corollary~1.2]{BaiXu2026}. Its homology is graded modulo twice the minimal Chern number, so classes in different integer degrees may enter the same group. The geometric stable Morse number retains the integer degrees. For our examples the two bounds differ: Remark~\ref{rem:periodic-bound} computes the cyclically graded integral bound and shows that our fixed-point count exceeds it.

A separate construction gives nondegenerate Hamiltonian counterexamples to the ordinary Morse bound in real dimension twelve, with arbitrarily large deficit \cite[Theorem~1.1 and Corollary~1.2]{OpenAI2026Morse}. Those examples use a nontrivial fundamental group and the distinction between ordinary and stable Morse theory. In the construction below, the ambient manifold and the four fixed components of an auxiliary involution are simply connected, and Smale's theorem identifies both Morse numbers with the minimum allowed by integral homology \cite[Theorem~2.4]{Smale1963}. The separation instead comes from distributing torsion at two primes across different integer degrees and realizing the smaller Morse counts on the fixed components of an involution. The proof below is independent of that companion construction.


\subsection*{The construction and its mechanism}
For a closed manifold $W$, let $b(W)$ be its total integral homology rank, and let $t_i(W)$ be the minimum number of generators of $\Tor H_i(W;\Z)$. The quantity
\[
 \lambda(W)=b(W)+2\sum_i t_i(W)
\]
is a lower bound for $\SM(W)$. For simply connected manifolds of the dimensions used here, Smale's theorem gives an ordinary Morse function attaining this bound. We prove the geometric lower bound in Section~\ref{sec:morse} and recall a handle proof of attainment in Appendix~\ref{app:handles}.

The dynamical ingredient is a Kähler surface $Y$, obtained by blowing up the four points $\{0,\infty\}\times\{0,\infty\}$ of $\CP^1\times\CP^1$. A half-turn of the diagonal circle action lifts to a Hamiltonian involution whose fixed set consists of four exceptional projective lines. For the high-dimensional simply connected Kähler manifold $X$ constructed below, the product involution $g$ on $M=X\times Y$ has four fixed components, each diffeomorphic to $C=X\times\CP^1$. Smale's theorem supplies a Morse function on $C$ with $\lambda(C)$ critical points.

Extend copies of this function to a $g$-invariant Hamiltonian on $M$, and let $\psi^s$ denote its flow. The flow commutes with $g$, so $(\psi^\epsilon\circ g)^2=\psi^{2\epsilon}$. For small $\epsilon$, the absence of short nonconstant periodic trajectories on the compact manifold $M$ forces every fixed point of $\psi^\epsilon\circ g$ to be an equilibrium on a fixed component of $g$. Section~\ref{sec:involution} proves that these are precisely the $4\lambda(C)$ chosen critical points and that they are all nondegenerate.

It remains to arrange $4\lambda(C)<\lambda(M)$. Section~\ref{sec:geometry} constructs $X$ by blowing up projective space along two disjoint centers with $2$-primary and $3$-primary torsion. The relevant elementary fact is
\[
 d\bigl((\Z/2)^a\oplus(\Z/3)^q\bigr)=\max(a,q),
\]
where $d$ denotes the minimum number of generators. Multiplication by $\CP^1$ and by $Y$ combines the torsion groups in nearby degrees with different multiplicities. When the two primes are combined into cyclic summands, these products permit different savings in the number of generators. Section~\ref{sec:counts} shows that only the ends of the degree range contribute to the resulting difference
\[
 \lambda(X\times Y)-4\lambda(X\times\CP^1)=16.
\]
The final assembly appears in Section~\ref{sec:assembly}. This use of mixed-prime torsion, together with the explicit involution, is the mechanism that puts the Hamiltonian count below the geometric stable Morse number. The proof does not require Floer theory.
